\begin{table}[htbp]
\centering
\small
\caption{Problema A, Passos 2 a 5 (google/gemini-2.5-flash-lite): métrica em R$^2$, R$^3$ e R$^4$ por regime.}
\label{tab:bloco1_A_passo4_R2_R3_R4_flashlite}
\begin{tabular}{llrrrrr}
\toprule
 &  & \multicolumn{2}{c}{Perceptron} & \multicolumn{2}{c}{NNLS} &  \\ \cmidrule(lr){3-4} \cmidrule(lr){5-6}
Espaço & Regime & Fid.\ (vs LLM) & Acur.\ vs real & Fid.\ (vs LLM) & Acur.\ vs real & $\Delta$ fid.\ vs R$^2$ \\
\midrule
R$^2$ & zero-shot† & 88,2 (87,3–88,7) & 14,2 (7,3–20,0) & 87,3 (86,0–88,7) & 10,7 (7,3–13,3) & +0,0 \\
R$^2$ & 4-shot† & 95,4 (85,6–100,0) & 93,2 (66,0–98,0) & 95,7 (82,9–100,0) & 93,2 (70,0–98,0) & +0,0 \\
R$^2$ & 10-shot† & 98,1 (94,3–100,0) & 96,0 (93,3–98,0) & 98,3 (93,6–100,0) & 95,7 (93,3–98,0) & +0,0 \\
R$^2$ & 20-shot† & 97,8 (94,6–100,0) & 96,8 (94,7–98,0) & 97,9 (93,8–100,0) & 96,6 (94,0–98,0) & +0,0 \\
R$^2$ & 40-shot† & 98,0 (94,5–100,0) & 96,9 (96,0–98,0) & 98,0 (94,5–100,0) & 96,9 (96,0–98,0) & +0,0 \\
R$^2$ & W ajustado ao rótulo real (teto)† & — & 97,3 (96,7–98,0) & — & 96,9 (96,7–97,3) & — \\
\midrule
R$^3$ & zero-shot† & 92,9 (91,3–94,7) & 18,0 (16,7–19,3) & 91,3 (90,7–92,0) & 15,6 (12,7–20,0) & +4,7 \\
R$^3$ & 4-shot† & 96,2 (91,1–100,0) & 93,0 (66,0–97,3) & 96,3 (88,4–100,0) & 93,5 (72,7–98,0) & +0,8 \\
R$^3$ & 10-shot† & 98,3 (95,7–100,0) & 95,2 (90,7–98,0) & 98,3 (94,3–100,0) & 95,6 (93,3–98,0) & +0,2 \\
R$^3$ & 20-shot† & 97,9 (95,4–99,2) & 96,3 (94,7–98,0) & 97,9 (93,8–100,0) & 96,6 (94,0–98,0) & +0,2 \\
R$^3$ & 40-shot† & 98,3 (97,3–100,0) & 95,9 (92,7–98,0) & 98,1 (95,5–100,0) & 96,9 (96,0–98,0) & +0,3 \\
R$^3$ & W ajustado ao rótulo real (teto)† & — & 97,3 (96,7–98,0) & — & 96,9 (96,7–97,3) & — \\
\midrule
R$^4$ & zero-shot† & 90,0 (87,3–93,3) & 13,3 (8,7–20,0) & 86,7 (84,7–88,7) & 10,4 (7,3–12,7) & +1,8 \\
R$^4$ & 4-shot† & 97,5 (93,2–99,3) & 91,0 (65,3–97,3) & 96,8 (91,1–100,0) & 91,6 (62,0–98,0) & +2,1 \\
R$^4$ & 10-shot† & 98,1 (94,3–100,0) & 95,3 (93,3–98,0) & 98,1 (93,6–100,0) & 95,6 (93,3–98,0) & +0,0 \\
R$^4$ & 20-shot† & 98,5 (96,2–100,0) & 95,8 (92,7–98,0) & 98,1 (96,2–100,0) & 95,9 (92,0–98,0) & +0,7 \\
R$^4$ & 40-shot† & 98,1 (95,5–100,0) & 95,9 (92,7–98,0) & 98,5 (96,4–100,0) & 96,4 (94,0–98,0) & +0,1 \\
R$^4$ & W ajustado ao rótulo real (teto)† & — & 97,3 (96,7–98,0) & — & 96,9 (96,7–97,3) & — \\
\bottomrule
\multicolumn{7}{p{0.95\linewidth}}{\footnotesize Problema A (linear). O LLM vê só x1, x2; R$^3$ acrescenta x1·x2 e R$^4$ acrescenta x1², x2² apenas na estimação. Em problema linear espera-se ganho ~nulo com a dimensão. Fid.\ (vs LLM) = fidelidade; Acur.\ vs real = acurácia (vs rótulo real) da métrica, contexto; o teto (W ajustado ao rótulo real) só tem acurácia. $\Delta$ = fidelidade Perceptron $-$ R$^2$. Passos do Problema A: zero-shot = Passo 2 (o LLM rotula sem exemplos); $k$-shot = Passo 3 (o LLM rotula com $k$ exemplos do rótulo real); a métrica de cada linha é estimada no Passo 4; teto = Passo 5 (W ajustado diretamente ao rótulo real). † = em pelo menos uma estimativa o Perceptron não encontrou $\hat{W}$ que separe os rótulos: melhor esforço ($\gamma = 0$) ou margem bruta mínima negativa (algum ponto do lado errado da fronteira).} \\
\end{tabular}
\end{table}
